\section{What this report establishes, and what it does not}
\label{sec:concl}

\textbf{Purpose.} This section states the whole result in one place, after the loads that bear on
it, so that no single section can be read as the conclusion on its own.

\subsection{What is established}
\label{sec:concl:established}

\begin{enumerate}
  \item \textbf{Drag: W is the cheapest geometry at every operating point, and the low-speed
    ordering carries to cruise.} At Condition~CR four candidates improve on the rigid baseline,
    W by $0.868$, F by $0.815$, C by $0.191$ and H by $0.009$ counts against $188.031$, and the
    whole field spans $1.1$ counts, enough to separate the groups but not the neighbours within
    them. At early cruise W saves $0.769$ counts against $177.839$, F and H cost $1.796$ and
    $2.244$, and C and M cost $12.871$ and $14.554$; late cruise reproduces it, W at $-0.781$ to
    M at $+13.728$ against $187.959$. \textbf{The two cruise points agree exactly}, both giving W,
    F, H, C, M, and they differ from the low-speed sequence W, F, C, H, M only in the C and H pair,
    which is $0.182$ counts apart at Condition~CR (Section~\ref{sec:perf:inversion}).
  \item \textbf{The cruise cost is the aft shock.} C and M end the supersonic region at
    $\eta = 0.80$ in a strong shock and separate behind it; F and H in a moderate shock and stay
    attached; W and the baseline without a shock. Drag follows that order exactly
    (Section~\ref{sec:perf:mechanism}).
  \item \textbf{Hinge moment: every candidate raises the actuator load, and two of the rows are
    not actuator loads at all.} All five raise the moment about the common $x_h/c = 0.62$ line at
    all three operating points, by $2.0$ to $5.0$ per cent at Condition~CR, $0.5$ to $6.3$ at
    early cruise and $0.4$ to $6.2$ at late cruise, with the same ordering at all three up to one
    near-tie. C, F and M act on the line integrated about, so their numbers are the
    moment their actuator reacts. \textbf{H's and W's are not}: H's hinge is at $0.70$, and W is
    not hinged at all but rotates the whole section about $0.25c$
    (Section~\ref{sec:hinge:hcaveat}).
  \item \textbf{Root bending: the constraint that selected C is not satisfied viscously.} C sits
    at $+6.796$ per cent in the VLM, inside the study's $6.8$ per cent screen by $0.004$
    percentage points, and the RANS puts it at $+7.170$ per cent, outside. M is outside on both.
    F, H and W are comfortably inside on both. At both cruise points, under the study's own
    stricter limit, all five exceed the baseline moment, W by the least
    (Section~\ref{sec:rbm:cruise}).
  \item \textbf{The two methods agree on root bending to better than half a point.} VLM and RANS
    differ by at most $0.40$ percentage points, so the screen result above is a property of the
    shapes rather than of the method.
  \item \textbf{The low-order method orders the candidates on its far-field column and prices
    none of them.} Far-field $C_{Diw}$ reproduces the RANS order at Condition~CR to nine pairs of
    ten, its one error on the closest pair in the field; near-field $C_{Di}$, the metric the June
    selection was made on, gives four of ten. Neither resolves the size of the effect: the
    far-field increments miss the measured ones by $86$ to $131$~per~cent of the whole low-speed
    span and the near-field by $63$ to $381$ (Section~\ref{sec:qv:decision}).
\end{enumerate}

\subsection{What is not established, stated as plainly}
\label{sec:concl:not}
\label{sec:perf:limits}

\begin{enumerate}
  \item \textbf{W's actuator load is not computed anywhere in this report.} Distributed twist
    reacts a torsional moment about $x/c = 0.25$ taken over the entire section. The quantity
    tabulated for W is about $0.62$ over the aft $38$ per cent, which is a valid comparison on a
    shared line and is not that load. Until the torsional moment is computed, \textbf{trailing-edge
    morphing and distributed twist cannot be ranked against one another on actuator load}, which is
    the axis the concept selection was expected to turn on.
  \item \textbf{H's actuator load is likewise not its tabulated number}, for the smaller reason
    that its own hinge is $0.08c$ aft of the comparison line.
  \item \textbf{Neither root-bending bound is a certified allowable.} The design study calls the
    $6.8$ per cent figure ``a study-level bound rather than a certified structural limit''. The
    cruise limit is a different and stricter definition. Which to size against is a question for
    the structures and actuator group and is not answered here.
  \item \textbf{The low-order comparison is made at Condition~CR only.} The VSPAERO decks and the
    meshed geometries share one frame and one geometry set there and nowhere else: the cruise
    decks carry the TP-1580 $S_\mathrm{ref}$ and a geometry set that maps to none of the meshed
    cases, and the study's cruise set contains optimiser seeds and the rigid baseline rather than
    the five candidates. Extending the comparison to either cruise point requires the
    corresponding VLM columns at that operating point, not merely the RANS ones, and quoting the
    two together without them would cross operating points (D076).
\item The meshed baseline is the uncorrected aerofoil-interpolation lineage and all five
candidates are the corrected one. Measured directly between the two \texttt{.vsp3} files
by \texttt{scripts/measure\_lineage\_pair.py} into \texttt{data/lineage\_measured.json},
the uncorrected wing carries 9 distinct aerofoils across its 19
stations where the corrected wing carries 19, so the two surfaces agree exactly at every
station they share and diverge only in between. They are identical at and inboard of
$\eta = 0.450$; the first station that differs is $\eta = 0.515$, which is also the worst,
at $0.407$ percentage points of $t/c$ and $0.209$ per cent of chord in surface ordinate. The
camber difference is an order of magnitude smaller again, at most $0.0034$ per cent of chord.
Both wings are the same planform at the same twist, so this is a thickness-channel
difference, and the candidate-to-candidate comparison is taken between surfaces of the
corrected lineage throughout and does not cross it.
\item Condition CR runs kOmegaSST wall resolved and both cruise points run Spalart--Allmaras wall
modelled, so counts are differenced within a condition only.
\item Every case is fully turbulent from the leading edge and in free air, so a tunnel comparison
needs blockage, wall-interference and support corrections on the measured side.
\item The VSPAERO columns are induced-drag increments and the RANS columns are total-drag
increments, so Table~\ref{tab:lo:selection} compares ordering rather than magnitude.
\item The cruise drag-due-to-lift slopes are formed on the trim legs alone, and six of the twelve
rest on a nudge too small to resolve the slope. They are used only to convert trim residuals of
at most $0.895$ counts of $C_L$ into a drag bias of at most $0.036$ counts.
\end{enumerate}

\subsection{Recommendation}
\label{sec:perf:recommendation}

\textbf{Recommendation.} W is the geometry to carry forward. It is the cheapest of the six at all
three operating points, $0.868$, $0.769$ and $0.781$ counts below the baseline at Condition~CR,
early cruise and late cruise, it raises $L/D$ at both cruise points, and it carries no aft shock
and no separation at cruise. It also has the smallest root-bending increase of the five at both
cruise points, $+3.36$ and $+3.25$ per cent, and sits inside the $6.8$ per cent screen at
Condition~CR on both methods. F and H are the next two, within $2.3$ counts of the baseline at
cruise. C and M are not candidates for cruise: they cost twelve to fifteen counts, separate behind
a strong aft shock, and both fail the root-bending screen once it is computed viscously.

\textbf{The caveat on W, and it is the one this report cannot close.} W's actuator reacts a
torsional moment about the quarter chord over the whole section, and that moment is not computed
here (Section~\ref{sec:hinge:hcaveat}); its tabulated hinge moment is a common-basis comparison
only. W's aerodynamic case is complete. Its mechanical case is not, and it is the input the
structures and actuator group needs before W can be committed to.

\subsection{The one number that would change the picture}
\label{sec:concl:next}
W's torsional moment about $x/c = 0.25$ over the full section, at every operating point. It is
the same integral this report already performs, over a different region about a different line, and
it is the missing input to the concept selection, and the recommendation above rests on it. Nothing
else outstanding has comparable leverage: the drag case is closed at all three operating points,
and the root-bending case is closed against both bounds.
